\documentclass[12pt]{article}
\usepackage[T1]{fontenc}
\usepackage[utf8]{inputenc}
\usepackage{lmodern}
\usepackage{textcomp}
\usepackage{enumitem}
\usepackage{geometry}
\geometry{margin=1in}
\begin{document}
\begin{center}
Answer Key - Political Science - Undergraduate Level Assessment 7 \
\small (Answers are ordered exactly as in the question paper.)
\end{center}

\section*{Multiple Choice}
\begin{enumerate}[label=\arabic*.]
\item \textbf{Answer:} C
\\ \textit{Options:} A) cooperation between the two major political parties; B) judicial activism; C) the constitutional fragmentation of power; D) affirmative action programs
\\ \textit{Note:} The constitutional fragmentation of power, through the separation of powers among the executive, legislative, and judicial branches, creates a system of checks and balances that limits the concentration of authority and fosters stability in American public policy by preventing any single entity from dominating the decision-making process.
\item \textbf{Answer:} D
\\ \textit{Options:} A) lack of a chief executive office; B) national government's inability to levy taxes effectively; C) absence of a central authority to regulate interstate trade; D) omission of a universal suffrage clause
\\ \textit{Note:} The omission of a universal suffrage clause was not addressed by the Constitution because the framers did not establish a national voting standard, allowing states to determine their own suffrage laws, which reflects the decentralized nature of power intended in the original document.
\item \textbf{Answer:} C
\\ \textit{Options:} A) People holding advanced academic degrees; B) Northeastern city dwellers; C) Southerners; D) Jewish Americans
\\ \textit{Note:} The Democratic Party's pursuit of liberal social policies often conflicts with the conservative values prevalent among Southerners, leading to their alienation from the party's traditional base.
\item \textbf{Answer:} C
\\ \textit{Options:} A) An increase in non-traditional security threats.; B) A more cooperative international system.; C) A more competitive international system.; D) An increase in inter-state alliances and military restraint.
\\ \textit{Note:} Inefficient balancing or buckpassing by states can lead to a more competitive international system because it creates power vacuums and uncertainty, prompting states to engage in more aggressive strategies to secure their interests and influence.
\item \textbf{Answer:} C
\\ \textit{Options:} A) President pro tempore; B) Speaker of the House; C) Congressional Management Foundation Chair; D) Minority Whip
\\ \textit{Note:} The Congressional Management Foundation Chair is not an official title or position held by a member of Congress, whereas special positions such as Speaker of the House or Senate Majority Leader are formal roles within the legislative branch.
\end{enumerate}

\section*{True / False}
\begin{enumerate}[label=\arabic*.]
\setcounter{enumi}{5}
\item \textbf{Answer:} False
\\ \textit{Options:} A) True; B) False
\\ \textit{Note:} The statement is false because the FBI primarily focuses on domestic law enforcement and counterterrorism efforts rather than on foreign intelligence, which is the responsibility of the Central Intelligence Agency (CIA).
\item \textbf{Answer:} True
\\ \textit{Options:} A) True; B) False
\\ \textit{Note:} Civil court is specifically designed to resolve disputes between private parties, such as individuals or organizations, rather than criminal cases brought by the state.
\item \textbf{Answer:} True
\\ \textit{Options:} A) True; B) False
\\ \textit{Note:} The First Amendment explicitly guarantees the right to petition the government, which allows individuals to express their concerns and seek remedies for grievances, thereby protecting this fundamental democratic principle.
\item \textbf{Answer:} False
\\ \textit{Options:} A) True; B) False
\\ \textit{Note:} Federal election laws do not prohibit negative advertising or 'attack' ads; instead, they regulate the disclosure of funding sources and require transparency regarding the sponsors of such advertisements.
\item \textbf{Answer:} False
\\ \textit{Options:} A) True; B) False
\\ \textit{Note:} The phrase 'surgical strike' is often criticized as a legitimate form of military engagement because it implies precision and minimal collateral damage, which can be misleading and does not account for the complex moral and legal implications of military interventions.
\end{enumerate}

\section*{Long Form Response}
\begin{enumerate}[label=\arabic*.]
\setcounter{enumi}{10}
\item \textbf{Answer:} Third-party movements play a significant role in the political landscape by introducing new ideas and priorities that can resonate with the electorate. These movements often emerge in response to perceived gaps in representation by the major parties, advocating for specific goals such as social justice, environmental issues, or economic reforms. When one or both of the major parties adopt these goals, it can lead to the disintegration of the third-party movement, as its original purpose is co-opted. This dynamic illustrates how third parties can influence the strategies of major parties, compelling them to adapt their platforms to capture the support of disenchanted voters. Ultimately, while third- party movements can stimulate political discourse and change, their survival is often contingent on their distinctiveness, which diminishes as their ideas are absorbed by larger parties.
\\ \textit{Note:} correct because it illustrates how third-party movements respond to unmet needs in the electorate, influence major parties to adopt their goals, and ultimately face disintegration when their core issues are co-opted, highlighting the dynamic interplay between third-party initiatives and the strategies of established political parties.
\item \textbf{Answer:} Energy security is intrinsically connected to various security sectors, including national security, economic security, and environmental security, each influencing one another in significant ways. National security is often contingent on access to reliable energy sources, as energy is essential for military operations and infrastructure. Economically, energy security impacts a nation's stability and growth; fluctuations in energy supply can lead to economic crises, which may in turn affect political stability. Additionally, environmental security is increasingly relevant, as energy production and consumption contribute to climate change, posing risks to both human security and geopolitical stability. These interconnections illustrate that energy security is not an isolated concern but a multifaceted issue that intertwines with broader security paradigms in political science.
\\ \textit{Note:} correct because it effectively illustrates the interdependence of energy security with national, economic, and environmental security, highlighting how reliable energy access is crucial for military readiness, economic stability, and environmental sustainability.
\end{enumerate}

\vfill
\noindent\textit{End of Answer Key}
\end{document}